List some factors contributing to the financial crisis 2007-2008

- Lower borrower quality (subprime mortgages)
- Shorter and shorter term financing of the banking sector
- Regulatory arbitrage
- More supply of financing subprime mortgages due to structured
products
- Top credit ratings for structured products based on subprime
loans

Vicious circle: lower prices, less collateral, more concerns
about solvency, and ever increasing haircuts

A 20\% haircut in a \$10 trillion market
means a shortage
of financing as
high as \$2 trillion.

What made the situation even riskier in 2008 than in earlier disruptions was the
rapid growth in the use of overnight repos
This movement toward shorter term financing meant that a greater portion of the
primary dealers funding was rolled over each day - risky development in a
market subject to the "repo run" dynamics. 

A second factor that made conditions in 2008 particularly precarious was the
resort to less liquid collateral in repo agreements

Asset Backed Commercial Papers (ABCP)
ABCP is typically a short-term instrument that matures between 1 and 270 days (average of 30 days),
is issued by an Asset-backed commercial paper program, such as a Conduit.
The financial assets that serve as collateral for ABCP are ordinarily a mix of many different assets, 
mostly asset-backed securities (ABS), CDO's. Most of the assets are AAA-rated.

Banks may be vulnerable to runs not based in fundamentals suggests
that ABCP programs may be vulnerable as well.


How do changes in interest rates affect banks balance sheet?

Banks assets
- typically long term assets
- Higher interest rates lead to lower values

Banks liabilities
- typically short term
- Higher interest rates lead to lower values
- low sensitivity to interest rate changes

Banks equity value (difference between assets and
liabilities) changes with interest rate if sensitivities of assets
and liabilities do not match due to maturity transformation.


How's duration defined in particular in the context of bonds.

\( \Delta P = P_{\text{post}} -  P_{\text{pre}}\)
\( \Delta r = r_{\text{post}} -  r_{\text{pre}}\)

\( P(r) = \sum_{t=1}^n C_t (1 + r)^{-t} \)
\( \partial_r P = \sum_{t=1}^n C_t (-t)(1 + r)^{-t-1} \)
\( - \partial_r P = \frac{1}{1+r} \sum_{t=1}^n t C_t (1 + r)^{-t} = D_{\text{abs}}(r)\)
absolute duration.
It holds approximately
\( \Delta P(r) \approx - D_{\text{abs}}(r) \delta r \)

(Macaulay) Duration
relative price change \( \frac{\Delta P(r)}{P(r)} \)
\( D = (\sum_{t=1}^n \frac{t C_t(1+r)^{-t}}{P(r)} = \sum_{t=1}^n t \frac{C_t(1+r)^{-t}}{\sum_{t=1}^n C_t(1+r)^{-t}} \)

Interpretation:
Weighted average of maturities (in years) with fraction of
present value of cash flows being the weights.

The duration of a portfolio can be calculated by taking the present
value weighted mean of the durations of the bonds in the portfolio.
Example:
Asset MarketValue Duration
Cash 10 0
AutoLoans 120 2
Mortgages 170 8
Liabilities MarketValue Duration
Checking\&Savings 120 0
CertificatesofDeposit 90 1
Long-Term-Financing 75 12

Duration of Assets:
\( \frac{10}{300}\cdot 0 + \frac{170}{300}\cdot 8 = 5.33 \)
Duration of Liabilities:
\( \frac{120}{285}\cdot 0 + \frac{90}{285}\cdot 1 + \frac{75}{285}\cdot 12 = 3.47 \)
Therefore, if market interest rate increases, assets will decrease
faster than liabilities.
It holds
\( \frac{\Delta P}{P} \approx - D \cdot \frac{\partial r}{1 + r} \)

\( \frac{E - 15}{15} \approx -40.67 \cdot \frac{0.055 - 0.05}{1.05} \implies E = 12.095 \)


Give an example of Immunization 

Duration of Liabilities:
\( \frac{120}{285}\cdot 0 + \frac{90}{285}\cdot 1 + \frac{75}{285}\cdot 12 = 3.47 \)

\( D_{\text{Equity}} = \frac{300}{300 - 285}\cdot D_{\text{Asset}} - \frac{285}{300-285}\cdot 3.47 = 0 \)
\( D_{\text{Assets}} = \frac{285\cdot 3.47}{300} = 3.3 \)

Let \( x \) be the PV of cash, PV \( y \) of mortgages.
\( \frac{x}{300}\cdot 0 + \frac{120}{300} \cdot 2 + \frac{y}{300}\cdot 8 = 3.3 \)
\( x + y = 180 \)
It follows that \( 170 - y \) of mortgage value should be transformed into cash (sold).

